% Options for packages loaded elsewhere
\PassOptionsToPackage{unicode}{hyperref}
\PassOptionsToPackage{hyphens}{url}
\documentclass[
]{article}
\usepackage{xcolor}
\usepackage[margin=1in]{geometry}
\usepackage{amsmath,amssymb}
\setcounter{secnumdepth}{-\maxdimen} % remove section numbering
\usepackage{iftex}
\ifPDFTeX
  \usepackage[T1]{fontenc}
  \usepackage[utf8]{inputenc}
  \usepackage{textcomp} % provide euro and other symbols
\else % if luatex or xetex
  \usepackage{unicode-math} % this also loads fontspec
  \defaultfontfeatures{Scale=MatchLowercase}
  \defaultfontfeatures[\rmfamily]{Ligatures=TeX,Scale=1}
\fi
\usepackage{lmodern}
\ifPDFTeX\else
  % xetex/luatex font selection
\fi
% Use upquote if available, for straight quotes in verbatim environments
\IfFileExists{upquote.sty}{\usepackage{upquote}}{}
\IfFileExists{microtype.sty}{% use microtype if available
  \usepackage[]{microtype}
  \UseMicrotypeSet[protrusion]{basicmath} % disable protrusion for tt fonts
}{}
\makeatletter
\@ifundefined{KOMAClassName}{% if non-KOMA class
  \IfFileExists{parskip.sty}{%
    \usepackage{parskip}
  }{% else
    \setlength{\parindent}{0pt}
    \setlength{\parskip}{6pt plus 2pt minus 1pt}}
}{% if KOMA class
  \KOMAoptions{parskip=half}}
\makeatother
\usepackage{color}
\usepackage{fancyvrb}
\newcommand{\VerbBar}{|}
\newcommand{\VERB}{\Verb[commandchars=\\\{\}]}
\DefineVerbatimEnvironment{Highlighting}{Verbatim}{commandchars=\\\{\}}
% Add ',fontsize=\small' for more characters per line
\usepackage{framed}
\definecolor{shadecolor}{RGB}{248,248,248}
\newenvironment{Shaded}{\begin{snugshade}}{\end{snugshade}}
\newcommand{\AlertTok}[1]{\textcolor[rgb]{0.94,0.16,0.16}{#1}}
\newcommand{\AnnotationTok}[1]{\textcolor[rgb]{0.56,0.35,0.01}{\textbf{\textit{#1}}}}
\newcommand{\AttributeTok}[1]{\textcolor[rgb]{0.13,0.29,0.53}{#1}}
\newcommand{\BaseNTok}[1]{\textcolor[rgb]{0.00,0.00,0.81}{#1}}
\newcommand{\BuiltInTok}[1]{#1}
\newcommand{\CharTok}[1]{\textcolor[rgb]{0.31,0.60,0.02}{#1}}
\newcommand{\CommentTok}[1]{\textcolor[rgb]{0.56,0.35,0.01}{\textit{#1}}}
\newcommand{\CommentVarTok}[1]{\textcolor[rgb]{0.56,0.35,0.01}{\textbf{\textit{#1}}}}
\newcommand{\ConstantTok}[1]{\textcolor[rgb]{0.56,0.35,0.01}{#1}}
\newcommand{\ControlFlowTok}[1]{\textcolor[rgb]{0.13,0.29,0.53}{\textbf{#1}}}
\newcommand{\DataTypeTok}[1]{\textcolor[rgb]{0.13,0.29,0.53}{#1}}
\newcommand{\DecValTok}[1]{\textcolor[rgb]{0.00,0.00,0.81}{#1}}
\newcommand{\DocumentationTok}[1]{\textcolor[rgb]{0.56,0.35,0.01}{\textbf{\textit{#1}}}}
\newcommand{\ErrorTok}[1]{\textcolor[rgb]{0.64,0.00,0.00}{\textbf{#1}}}
\newcommand{\ExtensionTok}[1]{#1}
\newcommand{\FloatTok}[1]{\textcolor[rgb]{0.00,0.00,0.81}{#1}}
\newcommand{\FunctionTok}[1]{\textcolor[rgb]{0.13,0.29,0.53}{\textbf{#1}}}
\newcommand{\ImportTok}[1]{#1}
\newcommand{\InformationTok}[1]{\textcolor[rgb]{0.56,0.35,0.01}{\textbf{\textit{#1}}}}
\newcommand{\KeywordTok}[1]{\textcolor[rgb]{0.13,0.29,0.53}{\textbf{#1}}}
\newcommand{\NormalTok}[1]{#1}
\newcommand{\OperatorTok}[1]{\textcolor[rgb]{0.81,0.36,0.00}{\textbf{#1}}}
\newcommand{\OtherTok}[1]{\textcolor[rgb]{0.56,0.35,0.01}{#1}}
\newcommand{\PreprocessorTok}[1]{\textcolor[rgb]{0.56,0.35,0.01}{\textit{#1}}}
\newcommand{\RegionMarkerTok}[1]{#1}
\newcommand{\SpecialCharTok}[1]{\textcolor[rgb]{0.81,0.36,0.00}{\textbf{#1}}}
\newcommand{\SpecialStringTok}[1]{\textcolor[rgb]{0.31,0.60,0.02}{#1}}
\newcommand{\StringTok}[1]{\textcolor[rgb]{0.31,0.60,0.02}{#1}}
\newcommand{\VariableTok}[1]{\textcolor[rgb]{0.00,0.00,0.00}{#1}}
\newcommand{\VerbatimStringTok}[1]{\textcolor[rgb]{0.31,0.60,0.02}{#1}}
\newcommand{\WarningTok}[1]{\textcolor[rgb]{0.56,0.35,0.01}{\textbf{\textit{#1}}}}
\usepackage{graphicx}
\makeatletter
\newsavebox\pandoc@box
\newcommand*\pandocbounded[1]{% scales image to fit in text height/width
  \sbox\pandoc@box{#1}%
  \Gscale@div\@tempa{\textheight}{\dimexpr\ht\pandoc@box+\dp\pandoc@box\relax}%
  \Gscale@div\@tempb{\linewidth}{\wd\pandoc@box}%
  \ifdim\@tempb\p@<\@tempa\p@\let\@tempa\@tempb\fi% select the smaller of both
  \ifdim\@tempa\p@<\p@\scalebox{\@tempa}{\usebox\pandoc@box}%
  \else\usebox{\pandoc@box}%
  \fi%
}
% Set default figure placement to htbp
\def\fps@figure{htbp}
\makeatother
\setlength{\emergencystretch}{3em} % prevent overfull lines
\providecommand{\tightlist}{%
  \setlength{\itemsep}{0pt}\setlength{\parskip}{0pt}}
\usepackage{bookmark}
\IfFileExists{xurl.sty}{\usepackage{xurl}}{} % add URL line breaks if available
\urlstyle{same}
% fallback for those not using the hyperref driver hyperxmp:
\makeatletter
\@ifundefined{xmpquote}{\newcommand{\xmpquote}[1]{#1}}{}
\makeatother
\hypersetup{
  pdftitle={Processing and Analysis of Biological Data},
  pdfauthor={Øystein H. Opedal, Department of Biology, Lund University},
  hidelinks,
  pdfcreator={LaTeX via pandoc}}

\title{Processing and Analysis of Biological Data}
\usepackage{etoolbox}
\makeatletter
\providecommand{\subtitle}[1]{% add subtitle to \maketitle
  \apptocmd{\@title}{\par {\large #1 \par}}{}{}
}
\makeatother
\subtitle{Chapter 3. The Linear Model II: Analysis of Variance}
\author{Øystein H. Opedal, Department of Biology, Lund University}
\date{2026-11-11}

\begin{document}
\maketitle

\subsection{Analysis of variance
(ANOVA)}\label{analysis-of-variance-anova}

When our predictor variables are categorical (factors), linear models
are used to perform analyses of variance. The parameter estimation works
in much the same way as in regression, except that instead of estimating
regression slopes, we are estimating group effects.

The aim of our ANOVA analysis is to partition the variance in our
response variable into a set of additive components (recall, variances
are additive). Based on this, we can evaluate whether the variance among
groups is greater than the variance within groups, or more so than
expected by chance.

The ANOVA framework is based on calculating ``sums of squares'' (SS),
the sum of the squared deviations of each group mean from the grand mean
(the variance among groups), and of each datapoint from either the grand
mean (which is equal to the total variance in the data), or the group
means (the residual or within-group variance).

\begin{figure}
\centering
\pandocbounded{\includegraphics[keepaspectratio,alt={Illustration of total, among-groups, and within-groups sums of squares. The total sum of squares is the total variance in the data, which can be partitioned into among-group and within-group (residual) components}]{Chapter3_ANOVA_extended_files/figure-latex/unnamed-chunk-1-1.pdf}}
\caption{Illustration of total, among-groups, and within-groups sums of
squares. The total sum of squares is the total variance in the data,
which can be partitioned into among-group and within-group (residual)
components}
\end{figure}

To understand how the ANOVA analysis works, let's simulate some data,
fit a linear model, and perform an ANOVA.

\begin{Shaded}
\begin{Highlighting}[]
\FunctionTok{set.seed}\NormalTok{(}\DecValTok{100}\NormalTok{)}
\NormalTok{groups }\OtherTok{=} \FunctionTok{as.factor}\NormalTok{(}\FunctionTok{rep}\NormalTok{(}\FunctionTok{c}\NormalTok{(}\StringTok{"Low"}\NormalTok{, }\StringTok{"Medium"}\NormalTok{, }\StringTok{"High"}\NormalTok{), }\AttributeTok{each=}\DecValTok{50}\NormalTok{))}
\NormalTok{x }\OtherTok{=} \FunctionTok{c}\NormalTok{(}\FunctionTok{rnorm}\NormalTok{(}\DecValTok{50}\NormalTok{, }\DecValTok{10}\NormalTok{, }\DecValTok{3}\NormalTok{), }\FunctionTok{rnorm}\NormalTok{(}\DecValTok{50}\NormalTok{, }\DecValTok{13}\NormalTok{, }\DecValTok{3}\NormalTok{), }\FunctionTok{rnorm}\NormalTok{(}\DecValTok{50}\NormalTok{, }\DecValTok{14}\NormalTok{, }\DecValTok{3}\NormalTok{))}

\FunctionTok{plot}\NormalTok{(groups, x, }\AttributeTok{las=}\DecValTok{1}\NormalTok{, }\AttributeTok{xlab=}\StringTok{""}\NormalTok{, }\AttributeTok{ylab=}\StringTok{"Body size (g)"}\NormalTok{)}
\end{Highlighting}
\end{Shaded}

\pandocbounded{\includegraphics[keepaspectratio]{Chapter3_ANOVA_extended_files/figure-latex/unnamed-chunk-2-1.pdf}}

Plots like this are called ``boxplots'', and can be useful for
visualising the distribution of data across factor levels or other
groups. With the default settings, the boxes range from the 1st to the
3rd quartile, i.e.~they span 50\% of the data. The thick lines show the
median, the ``whiskers'' extend to 1.5 times the inter-quantile range,
and individual circles show outliers. This representation allows us to
assess whether the data are roughly normally distributed within each
group, which will tend to yield normally distributed residuals within
each group, an assumption of the ANOVA model.

Boxplots are not directly useful for discussing ANOVA and variance
partitioning though, and plots that show all the data can be more
informative. Here is a rather elaborate plot combining a scatterplot
(with values slightly ``jittered'' along the x-axis for clarity) and a
boxplot.

EXTRA EXERCISE: Reproduce a plot similar to this for the ANOVA exercise.

\pandocbounded{\includegraphics[keepaspectratio]{Chapter3_ANOVA_extended_files/figure-latex/unnamed-chunk-3-1.pdf}}

Just as for linear regression, we fit the model with the \texttt{lm}
function and save the results to an object (here \texttt{m}). To perform
an ANOVA based on the fitted model, we call the \texttt{anova} function.

\begin{Shaded}
\begin{Highlighting}[]
\NormalTok{m }\OtherTok{=} \FunctionTok{lm}\NormalTok{(x}\SpecialCharTok{\textasciitilde{}}\NormalTok{groups)}
\FunctionTok{anova}\NormalTok{(m)}
\end{Highlighting}
\end{Shaded}

\begin{verbatim}
## Analysis of Variance Table
## 
## Response: x
##            Df  Sum Sq Mean Sq F value    Pr(>F)    
## groups      2  319.97 159.985  19.591 2.866e-08 ***
## Residuals 147 1200.43   8.166                      
## ---
## Signif. codes:  0 '***' 0.001 '**' 0.01 '*' 0.05 '.' 0.1 ' ' 1
\end{verbatim}

The ANOVA table contains a lot of information. First, we learn about the
number of degrees of freedom for each variable. For the \texttt{groups}
variable (our focal factor), the 2 degrees of freedom is the number of
groups in our data (3) - 1. The minus 1 comes from the fact that we had
to estimate the mean in the data to obtain our sums of squares (the sum
of the square deviations of data points from their group means).
Similarly for residual degrees of freedom, we have 150 - 2 - 1, where
the 2 comes from estimating the two contrasts (difference of group 2 and
3 from group 1), and the 1 is still the estimated mean.

The \texttt{Sum\ Sq} are the sums of squares, i.e.~the sum of the
squared deviations of each observation from the grand mean. The total
sum of squares (\(SS_T\)) divided by \(n-1\) gives the total variance of
the sample.

\begin{Shaded}
\begin{Highlighting}[]
\NormalTok{SS\_T }\OtherTok{=} \FloatTok{319.97+1200.43}
\NormalTok{SS\_T}\SpecialCharTok{/}\NormalTok{(}\DecValTok{150{-}1}\NormalTok{)}
\end{Highlighting}
\end{Shaded}

\begin{verbatim}
## [1] 10.20403
\end{verbatim}

\begin{Shaded}
\begin{Highlighting}[]
\FunctionTok{var}\NormalTok{(x)}
\end{Highlighting}
\end{Shaded}

\begin{verbatim}
## [1] 10.20403
\end{verbatim}

We can easily get the proportion of variance explained by the
\texttt{groups} variable, which is the same as the \(r^2\) for the
model.

\begin{Shaded}
\begin{Highlighting}[]
\FloatTok{319.97}\SpecialCharTok{/}\NormalTok{SS\_T}
\end{Highlighting}
\end{Shaded}

\begin{verbatim}
## [1] 0.2104512
\end{verbatim}

The \texttt{Mean\ Sq} is the variance attributable to each variable,
conventionally called the mean sum of squares (the sum of squares
divided by the degrees of freedom). The F-ratio (the test statistic in
an ANOVA) is computed as the mean sum of squares for the focal variable
divided by the mean residual sum of squares. Thus, it represents the
ratio of the among-group variance to the within-group variance, but also
the sample size which gives the residual degrees of freedom and thus,
all else being equal, larger sample gives lower mean sum of squares and
thus higher F-ratios and lower \(P\)-values.

In an ANOVA, a statistically supported among-group variance component
such as the one above indicates that at least one group mean is
different from the others. To further assess which groups are different,
we can extract the typical summary table of the linear model.

\begin{Shaded}
\begin{Highlighting}[]
\FunctionTok{summary}\NormalTok{(m)}
\end{Highlighting}
\end{Shaded}

\begin{verbatim}
## 
## Call:
## lm(formula = x ~ groups)
## 
## Residuals:
##     Min      1Q  Median      3Q     Max 
## -6.5887 -1.5596 -0.0987  1.6274  7.9729 
## 
## Coefficients:
##              Estimate Std. Error t value Pr(>|t|)    
## (Intercept)   13.7006     0.4041  33.901  < 2e-16 ***
## groupsLow     -3.4561     0.5715  -6.047 1.16e-08 ***
## groupsMedium  -0.9277     0.5715  -1.623    0.107    
## ---
## Signif. codes:  0 '***' 0.001 '**' 0.01 '*' 0.05 '.' 0.1 ' ' 1
## 
## Residual standard error: 2.858 on 147 degrees of freedom
## Multiple R-squared:  0.2105, Adjusted R-squared:  0.1997 
## F-statistic: 19.59 on 2 and 147 DF,  p-value: 2.866e-08
\end{verbatim}

This contains some of the same information as the ANOVA table, but we
now also obtain parameter estimates. The first parameter, the intercept,
corresponds to the estimated mean for the first level of the
\texttt{groups} factor. In this example this happens to be `High',
because H comes before L and M in the alphabet. The next two estimates
represents \emph{contrasts} from the reference group, and the associated
hypothesis tests test the null hypothesis that the group has the same
mean as the reference group.

The summary table also gives us directly the \(r^2\), which is a simple
ANOVA is defined as \(1-SS_E/SS_T\), where \(SS_E\) is the sum of
squares for the error term (residuals), and \(SS_T\) is the total sum of
squares. (Control question: why could we do it even more simply above?).

The parameter estimates allow us to quantify the effect size, i.e.~the
magnitude of the difference between the groups. A useful way to report
such differences is to compute the \% difference (the contrast divided
by the mean of the reference group, here -3.456/13.701 = -0.252), so
that we can say that ``Individuals in the low-food treatment were 25.2\%
smaller than individuals in the high-food treatment''.

Note that if we want a different reference group, we can change the
order of the factor levels.

\begin{Shaded}
\begin{Highlighting}[]
\NormalTok{groups }\OtherTok{=} \FunctionTok{factor}\NormalTok{(groups, }\AttributeTok{levels=}\FunctionTok{c}\NormalTok{(}\StringTok{"Low"}\NormalTok{, }\StringTok{"Medium"}\NormalTok{, }\StringTok{"High"}\NormalTok{))}
\NormalTok{m }\OtherTok{=} \FunctionTok{lm}\NormalTok{(x}\SpecialCharTok{\textasciitilde{}}\NormalTok{groups)}
\FunctionTok{summary}\NormalTok{(m)}
\end{Highlighting}
\end{Shaded}

\begin{verbatim}
## 
## Call:
## lm(formula = x ~ groups)
## 
## Residuals:
##     Min      1Q  Median      3Q     Max 
## -6.5887 -1.5596 -0.0987  1.6274  7.9729 
## 
## Coefficients:
##              Estimate Std. Error t value Pr(>|t|)    
## (Intercept)   10.2445     0.4041  25.349  < 2e-16 ***
## groupsMedium   2.5284     0.5715   4.424 1.88e-05 ***
## groupsHigh     3.4561     0.5715   6.047 1.16e-08 ***
## ---
## Signif. codes:  0 '***' 0.001 '**' 0.01 '*' 0.05 '.' 0.1 ' ' 1
## 
## Residual standard error: 2.858 on 147 degrees of freedom
## Multiple R-squared:  0.2105, Adjusted R-squared:  0.1997 
## F-statistic: 19.59 on 2 and 147 DF,  p-value: 2.866e-08
\end{verbatim}

Sometimes we also want to suppress the intercept of the model, and thus
estimate the mean and standard error for each level of the predictor. We
can do this by adding \texttt{-1} to the model formula (what comes after
the \texttt{\textasciitilde{}} sign). This could be useful for example,
if we wanted to obtain the estimated mean for each group, associated for
example with a 95\% confidence interval.

\begin{Shaded}
\begin{Highlighting}[]
\NormalTok{m }\OtherTok{=} \FunctionTok{lm}\NormalTok{(x}\SpecialCharTok{\textasciitilde{}}\NormalTok{groups}\DecValTok{{-}1}\NormalTok{)}
\FunctionTok{summary}\NormalTok{(m)}\SpecialCharTok{$}\NormalTok{coef}
\end{Highlighting}
\end{Shaded}

\begin{verbatim}
##              Estimate Std. Error  t value     Pr(>|t|)
## groupsLow    10.24453  0.4041333 25.34937 1.586298e-55
## groupsMedium 12.77295  0.4041333 31.60578 1.980446e-67
## groupsHigh   13.70064  0.4041333 33.90129 2.276884e-71
\end{verbatim}

\begin{Shaded}
\begin{Highlighting}[]
\FunctionTok{confint}\NormalTok{(m)}
\end{Highlighting}
\end{Shaded}

\begin{verbatim}
##                  2.5 %   97.5 %
## groupsLow     9.445865 11.04319
## groupsMedium 11.974287 13.57161
## groupsHigh   12.901979 14.49930
\end{verbatim}

\subsection{Two-way ANOVA}\label{two-way-anova}

Analyses of variance can also be performed with more than one factor
variable. If we have two factors, we can talk about two-way ANOVA, and
so on. A typical example from biology is when we have performed a
factorial experiment, and want to assess the effects of each
experimental factor and their potential interaction.

With two factors, a full model can be formulated as
\texttt{y\ \textasciitilde{}\ factor1\ *\ factor2}. Recall that in
\texttt{R} syntax, the * means both main effects and their interaction,
while a : means only the interaction. A detectable interaction term in
this model would indicate that the effect of factor 1 depends on the
level of factor 2 (and \emph{vice versa}). If we are analysing an
experiment where we have manipulated both temperature and nitrogen
supply, an interaction would mean that the effect of temperature depend
on the nitrogen level.

\subsubsection{Data exercise: analysing a factorial
experiment}\label{data-exercise-analysing-a-factorial-experiment}

The following data are from an experiment where female butterflies
reared on two different host plants were allowed to oviposit on the same
two host plants. The data include developmental time for the larvae, the
adult weight, and the growth rate of the larvae.

Analyse the data to assess effects of maternal vs.~larval host plant on
one or more response variable. Interpret the results and produce a nice
plot to illustrate.

\begin{Shaded}
\begin{Highlighting}[]
\NormalTok{dat }\OtherTok{=} \FunctionTok{read.csv}\NormalTok{(}\StringTok{"datasets/butterflies.csv"}\NormalTok{)}
\FunctionTok{names}\NormalTok{(dat)}
\end{Highlighting}
\end{Shaded}

\begin{verbatim}
## [1] "LarvalID"        "LarvalHost"      "Sex"             "MaternalHost"   
## [5] "MotherID"        "DevelopmentTime" "AdultWeight"     "GrowthRate"
\end{verbatim}

As a first step, let us compute some summary statistics, like the mean
development time for each combination of larval and maternal host plant.

\begin{Shaded}
\begin{Highlighting}[]
\NormalTok{dat}\SpecialCharTok{$}\NormalTok{MaternalHost }\OtherTok{=} \FunctionTok{paste0}\NormalTok{(dat}\SpecialCharTok{$}\NormalTok{MaternalHost, }\StringTok{"M"}\NormalTok{)}
\NormalTok{dat}\SpecialCharTok{$}\NormalTok{LarvalHost }\OtherTok{=} \FunctionTok{paste0}\NormalTok{(dat}\SpecialCharTok{$}\NormalTok{LarvalHost, }\StringTok{"L"}\NormalTok{)}
\NormalTok{means }\OtherTok{=} \FunctionTok{tapply}\NormalTok{(dat}\SpecialCharTok{$}\NormalTok{DevelopmentTime, }\FunctionTok{list}\NormalTok{(dat}\SpecialCharTok{$}\NormalTok{MaternalHost, dat}\SpecialCharTok{$}\NormalTok{LarvalHost), mean)}
\NormalTok{means}
\end{Highlighting}
\end{Shaded}

\begin{verbatim}
##           BarbareaL BerteroaL
## BarbareaM  21.69608  27.00000
## BerteroaM  23.51282  31.01923
\end{verbatim}

\begin{figure}
\centering
\pandocbounded{\includegraphics[keepaspectratio,alt={Larval developmental time depending on larval and maternal host plant}]{Chapter3_ANOVA_extended_files/figure-latex/unnamed-chunk-13-1.pdf}}
\caption{Larval developmental time depending on larval and maternal host
plant}
\end{figure}

To get started with the formal analyses, fit a linear model with larval
and maternal host, as well as their interaction, as predictors. Produce
an ANOVA table.

\begin{Shaded}
\begin{Highlighting}[]
\FunctionTok{names}\NormalTok{(dat)}
\end{Highlighting}
\end{Shaded}

\begin{verbatim}
## [1] "LarvalID"        "LarvalHost"      "Sex"             "MaternalHost"   
## [5] "MotherID"        "DevelopmentTime" "AdultWeight"     "GrowthRate"
\end{verbatim}

\begin{Shaded}
\begin{Highlighting}[]
\NormalTok{m }\OtherTok{=} \FunctionTok{lm}\NormalTok{(DevelopmentTime}\SpecialCharTok{\textasciitilde{}}\NormalTok{MaternalHost}\SpecialCharTok{*}\NormalTok{LarvalHost, }\AttributeTok{data=}\NormalTok{dat)}
\FunctionTok{anova}\NormalTok{(m)}
\end{Highlighting}
\end{Shaded}

\begin{verbatim}
## Analysis of Variance Table
## 
## Response: DevelopmentTime
##                          Df  Sum Sq Mean Sq F value    Pr(>F)    
## MaternalHost              1  623.61  623.61  177.90 < 2.2e-16 ***
## LarvalHost                1 2682.41 2682.41  765.21 < 2.2e-16 ***
## MaternalHost:LarvalHost   1   80.80   80.80   23.05 2.561e-06 ***
## Residuals               283  992.05    3.51                      
## ---
## Signif. codes:  0 '***' 0.001 '**' 0.01 '*' 0.05 '.' 0.1 ' ' 1
\end{verbatim}

To interpret the results, we look first at the ANOVA table. There are
detectable effects of both larval and maternal host plants, with more
variance explained by larval host than maternal host (based on the
larger sum of squares for larval host). Furthermore, the effects of
larval and maternal host are not independent, as indicated by a
detectable interaction. The variance explained by the interaction term
is limited though, as we can also see from the graph.

\newpage

\begin{Shaded}
\begin{Highlighting}[]
\FunctionTok{summary}\NormalTok{(m)}
\end{Highlighting}
\end{Shaded}

\begin{verbatim}
## 
## Call:
## lm(formula = DevelopmentTime ~ MaternalHost * LarvalHost, data = dat)
## 
## Residuals:
##     Min      1Q  Median      3Q     Max 
## -5.0192 -1.5128 -0.5128  1.0000  7.0000 
## 
## Coefficients:
##                                           Estimate Std. Error t value Pr(>|t|)
## (Intercept)                                21.6961     0.1854 117.033  < 2e-16
## MaternalHostBerteroaM                       1.8167     0.2816   6.451 4.80e-10
## LarvalHostBerteroaL                         5.3039     0.3132  16.934  < 2e-16
## MaternalHostBerteroaM:LarvalHostBerteroaL   2.2025     0.4588   4.801 2.56e-06
##                                              
## (Intercept)                               ***
## MaternalHostBerteroaM                     ***
## LarvalHostBerteroaL                       ***
## MaternalHostBerteroaM:LarvalHostBerteroaL ***
## ---
## Signif. codes:  0 '***' 0.001 '**' 0.01 '*' 0.05 '.' 0.1 ' ' 1
## 
## Residual standard error: 1.872 on 283 degrees of freedom
## Multiple R-squared:  0.7734, Adjusted R-squared:  0.771 
## F-statistic: 322.1 on 3 and 283 DF,  p-value: < 2.2e-16
\end{verbatim}

To interpret the summary table from a two-way ANOVA, we need to know how
to read the parameter estimates. The Intercept is, as before, the
estimated mean for the reference group. Here, the reference group is the
combination of the the first level of each factor, in this case
\emph{Barbarea} as both maternal and larval host plant. The following to
terms are the contrasts between the reference group and the alternative
maternal and larval host plants, respectively. Finally, the interaction
effect (the last term above), gives the \emph{added} effect of changing
both the maternal and the larval host plant. Hence, to obtain the
estimated mean for the combination maternal host \emph{Berteroa} and
larval host \emph{Berteroa}, we sum all the parameter estimates. Check
that this is indeed the case!

To quantify the effect size, it is useful to first compute the mean
development time for each larval and maternal host plant.

\begin{verbatim}
## BarbareaL BerteroaL 
##  22.60445  29.00962
\end{verbatim}

\begin{verbatim}
## BarbareaM BerteroaM 
##  24.34804  27.26603
\end{verbatim}

Based on this, we could write the methods and results like this:

\emph{Methods} To assess differences in development time between larvae
grown on \emph{Barbarea} and \emph{Berteroa}, and between larvae whose
mothers were grown on the same two hosts, we fitted a linear model with
development time as response variable, and larval and maternal hosts as
predictors, and performed an analysis of variance based on the fitted
model. To assess transgenerational effects, we also includes the
interaction term between maternal and larval host. Thus, in R syntax,
our model took the form

\[Development Time \sim Larval Host * Maternal Host\]

\emph{Results} The larvae developed 22.1\% faster when grown on Barbarea
than when grown on Berteroa (mean development time = 22.6 and 29.0 days,
respectively, \(F_{1,283}=765.21\), Fig. 1). Larvae whose mothers were
grown on Barbarea developed 10.7\% faster (mean development time = 24.3
and 27.3 days, respectively, \(F_{1,283}=177.90\)). The difference in
development time between larval host plants was slightly larger when the
mother was grown on Berteroa than when the mother was grown on Barbarea
(24.2\% vs.~19.6\% reduction in developmental time on Barbarea,
respectively).

\end{document}
